
\documentclass{beamer}
\usepackage{tikz}

\begin{document}

\begin{frame}
\frametitle{Inertia in Physics}
\begin{itemize}
  \item Inertia is a fundamental property of matter described by Sir Isaac Newton.
  \item It refers to the tendency of an object to resist changes in its state of motion.
\end{itemize}

\begin{tikzpicture}[remember picture, overlay]
  \draw[blue, ultra thick] (1,1) rectangle (5,3);
\end{tikzpicture}
\end{frame}

\begin{frame}
\frametitle{Definition of Inertia}
\begin{itemize}
  \item Inertia is the property of matter due to which an object resists any change in its state of rest or motion.
  \item An object will remain at rest or continue moving in a straight line at a constant velocity unless acted upon by an external force.
\end{itemize}

\begin{tikzpicture}[remember picture, overlay]
  \draw[red, ultra thick] (1,1) circle (1.5);
\end{tikzpicture}
\end{frame}

\begin{frame}
\frametitle{Effect of Mass on Inertia}
\begin{itemize}
  \item The greater the mass of an object, the greater its inertia.
  \item More massive objects are more resistant to changes in their motion.
\end{itemize}

\begin{tikzpicture}[remember picture, overlay]
  \draw[green, ultra thick] (2,2) ellipse (2 and 1);
\end{tikzpicture}
\end{frame}

\end{document}
